\textbf{*NaR-} & nṛ & naḻ & nar & [Man/Hero]. The stable alveolar nasal onset denotes the specific identity of the agential human. Vedic \textit{nṛ} preserves the vocalic liquid. \\ \addlinespace[12pt]

\textbf{*MoR-} & maru & māl & mar & [Mountain/Dark]. The bilabial nasal governs the concept of shadowy or massive topographical features. Stable sonorant final. \\ \addlinespace[12pt]

\textbf{*NiL-} & nil & nil & nyal & [Blue/Deep]. Deictic color cluster. The front vowel triggers palatalization in the Eastern node while remaining stable in the South. \\ \addlinespace[12pt]

\textbf{*K'eH-} & gā- & kā- & khyer- & [Give/Call]. The velar ejective transformation to Vedic \textit{g}. Laryngeal scarring produces length in the West and South (Vedic \textit{gā}, Tamil \textit{kāஃ}). \\ \addlinespace[12pt]

\textbf{*T'eH-} & dah- & tā- & thos- & [Burn/Fire]. Dental ejective radiation. The high-energy laryngeal coda triggers compensatory vowel lengthening in Old Tamil \textit{tāஃ}. \\ \addlinespace[12pt]

\textbf{*P'eR-} & bhar & peṟu & phar & [Bear/Result]. Labial ejective shift. This root identifies the resultative state of the agential 'carrying' action. \\ \addlinespace[12pt]

\textbf{*S-N-} & sūnu & maṉai & sun & [Son/Child]. Fricative-nasal cluster governing domestic outcome. Old Tamil \textit{maṉai} exhibits a common secondary nasal shift. \\ \addlinespace[12pt]

\textbf{*K'V-} & go- & ko- & khye- & [Go/Path]. Minimalist velar ejective root. The Western node preserves the high-volume voicing characteristic of the PED motion verbs. \\ \addlinespace[12pt]

\textbf{*T'V-} & de- & te- & the- & [This/Proximate]. The primary dental deictic. This root provides the foundation for the entire PED spatial topography system. \\ \addlinespace[12pt]

\textbf{*H-eN-} & anya & ēṉ & -an & [Other/Different]. Laryngeal-nasal cluster. The Western node exhibits the 'vowel coloring' effect typical of the PIE laryngeals. \\ \addlinespace[12pt]

\textbf{*K'uS-} & guṣ- & kuci & khus & [Grasp/Tight]. Velar-fricative cluster. The semantic core relates to the physical act of constriction or binding. \\ \addlinespace[12pt]

\textbf{*P'iL-} & bil- & pil & phel & [Split/Hole]. Labial ejective radiation. This root anchors the PED partitive system and the concept of 'void'. \\ \addlinespace[12pt]

\textbf{*T'uM-} & dum & tum & thum & [Deep/Dark]. Dental ejective governing the concept of obscurity. Note the perfect preservation of the labial nasal coda. \\ \addlinespace[12pt]

\textbf{*S-R-} & sar-at & car & ser & [Year/Cycle]. Fricative-liquid cluster. The Southern node demonstrates the car/sar palatalization rule established in Part II. \\ \addlinespace[12pt]

\textbf{*M-K-} & muc- & mukku & mkhye- & [Release/Free]. Nasal-velar cluster. The Western node exhibits a secondary palatalization shift common in the PIE 'Satem' group. \\ \addlinespace[12pt]

\textbf{*N-T'-} & nad-ī & nattu & nath- & [River/Flow]. Alveolar-dental cluster. The dental ejective surfaces as the voiced stop \textit{d} in the Vedic hydronyms. \\ \addlinespace[12pt]

\textbf{*K'eN-} & kṛ-t & kaṇ-tu & mkhyen- & [Know/Done]. Velar ejective transformation. The semantic links the 'completed action' to the 'acquired knowledge'. \\ \addlinespace[12pt]

\textbf{*P'uL-} & bhu- & pul-u & phul- & [Abundant/Full]. Labial ejective shift. The high-aspiration energy in Tibetan \textit{phul} identifies the original glottalic pressure. \\ \addlinespace[12pt]

\textbf{*S-K'-} & sag- & cak-u & sakh- & [Follow/Partner]. Sibilant-velar cluster. The ejective transformation follows the standard Glottalic Law across the three nodes. \\ \addlinespace[12pt]

\textbf{*M-T'-} & mad-hu & matt-u & meth- & [Sweet/Intoxicant]. The famous 'Honey' root. The dental ejective produces the characteristic Vedic \textit{d} and Tamil \textit{tt} geminate. \\ \addlinespace[12pt]

\textbf{*T'aL-} & dal-a & tal-ir & thal- & [Green/Leaf]. Dental ejective radiation. The semantic core identifies the 'outgrowth' or 'extension' of the botanical system. \\ \addlinespace[12pt]

\textbf{*K'aR-} & gar-j & kar & khar-u & [Cry/Roar]. Velar ejective mass cluster. The root governs high-volume vocalizations across the Altai-Dravidian continuum. \\ \addlinespace[12pt]

\textbf{*P'eN-} & bandh- & piṇ-i & phan- & [Bind/Tie]. Labial ejective radiation. Note the aspiration in Tibetan and the secondary voicing in the Western node. \\ \addlinespace[12pt]

\textbf{*S-u-} & sū- & cu- & su- & [Birth/Source]. Fricative-vowel root. The sibilant onset is preserved with high fidelity across 30,000 years of drift. \\ \addlinespace[12pt]

\textbf{*N-e-} & na- & na- & na- & [Identity/Exist]. The minimalist nasal root. The alveolar nasal remains the primary marker of 'presence' in the PED substrate. \\ \addlinespace[12pt]
